% LSPV-4 REPORT::SELECTOR-CAPACITY-BIT-LEDGER81::2026-10-24T15:36:00+09:00
\documentclass[11pt]{article}
\usepackage[margin=.86in]{geometry}
\usepackage{amsmath,amssymb,amsthm,booktabs,url}
\setlength{\emergencystretch}{3em}
\newtheorem{theorem}{Theorem}
\newtheorem{proposition}{Proposition}
\title{SELECTOR-CAPACITY-BIT-LEDGER81:\\Full-Child Descriptors Versus One-Bit SAT Decisions}
\author{LSPV-4: Transport, Endpoints, and Bit-Resource Accounting}
\date{24 October 2026}
\begin{document}\maketitle
\noindent\textbf{P versus NP remains OPEN.} The results below are restricted descriptor and source-aware decision theorems, not all-input ARC successors. Independent Role~5 review is pending.
\section{Claim card and evidence sources}
We audit whether bounded child serializers or charged source selectors can supply the missing favorable ARC/FLMV achievable-region edge. Three falsifiable possibilities are: (A) a stateless full-service decoder represents every $K$-subset of a fixed $M$-clause parent in $o(K\log(M/K))$ bits; (B) a large selector descriptor implies a difficult one-bit child SAT label; and (C) clause-deletion children reconstruct parent SAT from their separate labels. (A) is impossible within the exact fixed-parent descriptor interface; (B) and (C) fail on explicit SAT/UNSAT controls. The information-theoretic lower bound for arbitrary clause selectors was \emph{already present} in SELECTOR-VIRTUAL-CHILD62, not newly discovered here. New in this report are the exact both-label family, count formula, service-separation proof, bit-address audit, and BIT-RESOURCE1 accounting.
\section{Same-parent, two-label family}
For $M\ge8$, let $P_M$ have the eight proper ternary clauses of every sign pattern on $(x_1,x_2,x_3)$ and $M-8$ positive ternaries, each on its own disjoint triple of fresh variables. All clauses are distinct. Use $n=3+3(M-8)$ named variables, including unused variables in child headers. For any sorted $K$-subset $S\subseteq\{0,\ldots,M-1\}$, child $C_S$ consists exactly of the selected clauses with the same $n$-variable header. Let $r(S)=|S\cap\{0,\ldots,7\}|$.
\begin{theorem}[Exact same-parent child labels, PROVED/EXACT]
For all selectors $S$,
\[
 \#SAT(C_S)=(8-r(S))\,7^{K-r(S)}8^{(M-8)-(K-r(S))}.
\]
The child is UNSAT if and only if the eight core indices are all selected. Therefore for $K<8$ all $\binom MK$ children are SAT despite $P_M$ being UNSAT; for $K\ge8$ exactly $\binom{M-8}{K-8}$ children are UNSAT.
\end{theorem}
\begin{proof}
Each distinct core sign pattern forbids exactly one of eight assignments to the core triple. The $r$ selected patterns leave $8-r$ assignments. Each selected independent positive filler has seven satisfying assignments and each unselected filler eight. Multiply these independent counts.
\end{proof}
\section{Descriptor capacity and source-service quantifiers}
Fix $P_M$. A full-child descriptor is a selector-dependent finite binary string $d_S$ for a single \emph{fixed} deterministic decoder $g$ such that $g(P_M,d_S,i)$ returns every bit $i$ of the ordinary, self-contained serialized child $C_S$. It has no other selector-dependent state, nonuniform code, free oracle, or writable external store. Every $K$-subset must be served. Descriptor storage is charged writable bits.
\begin{theorem}[Universal selector capacity, PROVED/EXACT restricted]
If descriptor length is fixed at $D$ bits, then
\[
D\ge\left\lceil\log_2\binom MK\right\rceil
\ge K\log_2(M/K).
\]
The first bound is attainable by colexicographic rank encodings, although arbitrary-bit unranking is not automatically polylogarithmic. A sorted index array using $K\lceil\log_2 M\rceil$ bits does support polylogarithmic source-backed bit access with $O(\log N)$ scratch.
\end{theorem}
\begin{proof}
Distinct selectors yield different ordered child clause lists, since the parent clauses are unique. Equal descriptors would produce identical bits under the deterministic decoder, a contradiction. Thus $\binom MK$ descriptors are different. The product identity $\binom MK=\prod_{i=0}^{K-1}(M-i)/(K-i)\ge(M/K)^K$ gives the second inequality. Combinatorial ranking $R(S)=\sum_{j=1}^K\binom{s_j}{j}$ for sorted zero-based indices is bijective. For fast random access, a requested child word identifies one array entry $s_j$, then its source clause word; all offsets and addresses have logarithmic width.
\end{proof}
This bound is \emph{not} universal for an algorithm that chooses only one selector, a codebook of $Q\ll\binom MK$ selectors, a streamed selector, or a decider that never reconstructs child bits.
\begin{proposition}[Cheap label-only service, PROVED/EXACT restricted]
Given $S$ as a sequential stream of selected indices, $SAT(C_S)$ is decided by one fixed uniform procedure in $O(K\log M)$ elementary bit work and $O(\log M)$ live work bits. No child CNF, colex rank, or full descriptor is needed.
\end{proposition}
\begin{proof}
Track an eight-bit mask indicating which core clause indices appeared. Output UNSAT iff all eight flags are set. Each index costs $O(\log M)$ bit operations to receive and compare. Both labels are handled without an answer oracle.
\end{proof}
This is an explicit separation between \emph{reconstructing a complete child} and \emph{deciding that child}. The source-aware shortcut defeats any inference from descriptor entropy to mandatory SAT computation. The full parent is also directly recognizable as UNSAT by its core, with template recognition charged on ordinary inputs.
\section{Ordinary bits, exponent transport, and BIT-RESOURCE1}
Use four $w$-bit words per clause plus two header words and $w(n,m)=\lceil\log_2(64(n+m+1))\rceil$. Parent and child lengths are
\[
N=(4M+2)w(n,M),\qquad L=(4K+2)w(n,K).
\]
Literal codes are $2|v|+[v<0]$ and each clause record is $(3,\ell_1,\ell_2,\ell_3)$. Distinct selectors produce distinct ordinary child bitstrings of the same length. With $K=\Theta(N^\rho/\log N)$, $M=\Theta(N/\log N)$ and any fixed $0<\rho<1$,
\[
\log_2\binom MK=\Theta(N^\rho),\quad
K\lceil\log_2 M\rceil=\Theta(N^\rho),\quad
L=N^{\rho+o(1)}.
\]
At $\rho=1/2$, the $N^{1/2+o(1)}$ selector-descriptor buffer of SELECTOR-VIRTUAL-CHILD62 is capacity-tight in exponent \emph{for the all-subsets full-child service}. It is not mandatory for a semantic SAT decider.
Under a \emph{separately proved} all-input correct parent contractor with $Q=N^{3/5+o(1)}$ sequential children of length $N^{1/2+o(1)}$, explicit writing of a fresh sorted descriptor for each yields the conditional upper ledger
\[
t'\le\max\{r,11/10,3/5+t/2\},\qquad
s'\le\max\{\eta,\kappa,1/2,s/2\},
\]
where $r$ charges parent-local source acquisition, selector discovery, failures, gluing and reconstruction, $\eta$ generator scratch, and $\kappa$ ancestors/caches. The $11/10$ term is not a universal lower bound on the work of correlated selector sequences. Every child must be an ordinary self-contained input or have a charged exact virtual-bit evaluator. No available theorem here supplies the missing all-input parent-label gluing.
The declared elementary BIT-RESOURCE1 vector separately counts primitive work, input/workspace reads and writes, actual value-changing flips, peak live writable bits, cumulative bit-time, and program/representation size. In the specified eager serializer every one of $L$ target bits is written even when unchanged from blank, while only the one-bits actually flip. A lazy decoder or source-aware label policy is allowed to avoid those writes. If $D$ contiguous descriptor bits remain physically live for $H$ steps, their conditional cumulative bit-time contribution is $DH$; discarding or recomputing them changes the charged trajectory.
\section{Finite audit and adaptive-state ledger}
The standalone Python audit used the exact parent core, a separate model counter on $M=10$ controls, complete colex rank inversion, and ordinary bit-by-bit comparison with a stateless virtual child evaluator. It enumerated all $K$-subsets in each row:
\[
\begin{array}{rrrrrr}
M&K&\mathrm{SAT}&\mathrm{UNSAT}&D_{\min}&K\lceil\log_2 M\rceil\\
10&7&120&0&7&28\\
10&8&44&1&6&32\\
16&7&11440&0&14&28\\
16&8&12869&1&14&32\\
16&9&11432&8&14&36
\end{array}
\]
Total \textbf{273719 assertion-accounted checks}, including \textbf{129030 exact virtual-bit comparisons}; zero failures (COMPUTATIONAL). Deterministic exhaustive enumeration, no random seed. Script SHA-256 \texttt{0881bc1602224a102c91b2af06cca02cd4eaef81004d6a058b13ca82dbdaa6a6}; raw JSON SHA-256 \texttt{f00daa10ef294fe744e7dc59920d5b9b12d753c21f867645827a28b7c78fba8a}. The full audit is in the accompanying local package; a compact executable also appears beside this report on GitHub.
\par\noindent\textbf{Adaptive ledger.} At the initial state the solver knows the claimed service, source pointer, $M,K$ and no SAT label. It charges a source/header scan and selector-index acquisition. If a one-bit child label is requested it chooses the eight-mask streaming rule and terminates after verification. If all child bits are requested, it chooses index-array virtual reads, rank unranking, or eager serialization, charging writes, unchanged writes, repeated source reads and peak cache. SAT controls omit one core clause; UNSAT controls contain all eight. This finite menu is an \emph{experimental restriction}, not coverage of every adaptive SAT algorithm; SCC, Horn, affine, learning, caching, and direct decision remain legitimate alternatives.
\section{Corpus review, gap acceptance, and next test}
Full/targeted originals read: root \texttt{main13.tex} (blob 87dbaad5e18932b1db7fca57fc40e0ac020a6f49), AGENTS, CODEX and research skill, SELECTOR-VIRTUAL-CHILD62 (blob c3303e3d3ac57ea484e1018a7d27dc59c0f07d09), VIRTUAL-CHILD-COMPOSITION61, CACHE-MISS-FRONTIER18, ROLE 4 SENSITIVITY-SATURATION80, ROLE 5 WRITE-SHADOW-ENDPOINT118, ROLE 3 FORBIDDEN-COMPUTATION84, ROLE 1 ARC-CHANNEL-CURVE119, ROLE 2 FLMV-CURVE-FRONTIER115 and the September 27 Drive ARC origin. Root \texttt{main.tex} was inspected by relevant sections, not fully manually reviewed. Earlier 3532-file crawl is an inherited keyword/SHA baseline, not a new exhaustive manual review. Drive LOG/BOOL/SIL pagination and historical main2--main12 remain unread in this targeted run.
\par\noindent\textbf{Duplication audit.} The entropy inequality and selector-array decoder are already in SELECTOR-VIRTUAL-CHILD62. This report adds the explicit common-parent SAT/UNSAT witness family, full model-count theorem, saturation of exact child-bit information, and a sharp label-only service bypass. CP22LIFT1 transports designated Cutting Planes proofs, not arbitrary SAT runtimes. The old $2\cos(\pi/7)$ universal product endpoint is retired. WRITE-SHADOW-ENDPOINT118 is a model-matched FLMV pullback, not an ARC successor.
\par\noindent\textbf{Gap ledger.} \textbf{FILLED, known:} FLMV Corollary 5.8 and conditional cofinal finite-hit arithmetic. \textbf{PROVED/EXACT restricted:} selector capacity, same-parent SAT/UNSAT counts, cheap status computation. \textbf{PARTIAL/RESTRICTED:} charged descriptor and caching simulations. \textbf{FALSIFIED/RETIRED:} descriptor capacity or separate short-child satisfiability as mandatory arbitrary-decider runtime/gluing. \textbf{OPEN:} existence of a uniform all-input, both-label, charged favorable ARC edge.
\par\noindent\textbf{Next artifact:} the full original HEREDITARY-SCOPED10 generator. Test on connected high-width both-label ordinary CNFs with explicit child bit service, paid source discovery, exact parent-label reconstruction, adaptive caching, and simultaneous peak space. The FLMV reference is Fortnow, Lipton, van Melkebeek and Viglas, \emph{Time-Space Lower Bounds for Satisfiability}, JACM 52(6) (2005), Corollary 5.8, in its random-access TM model, with $\beta(c)=((c+2)/c^2-c)/2$ and a strict endpoint box near $(6/5,1/2)$.
\begin{thebibliography}{4}
\bibitem{flmv} Fortnow, Lipton, van Melkebeek, Viglas, \emph{Time-Space Lower Bounds for Satisfiability}, JACM 52(6) (2005), 835--865. \url{https://doi.org/10.1145/1101821.1101822}.
\bibitem{vc62} Project original, SELECTOR-VIRTUAL-CHILD62. \url{https://github.com/j15marti/pnp/blob/main/research/main/2026-10-22_1048_MAIN_SELECTOR_VIRTUAL_CHILD62.md}.
\bibitem{vc61} Project original, VIRTUAL-CHILD-COMPOSITION61. \url{https://github.com/j15marti/pnp/blob/main/research/main/2026-10-22_0948_MAIN_VIRTUAL_CHILD_COMPOSITION61.md}.
\bibitem{main13} Project original, \emph{Main 13: ARC/FLMV Proof Chain}. \url{https://github.com/j15marti/pnp/blob/main/main13.tex}.
\end{thebibliography}
\end{document}
% END_OF_VERIFIED_REPORT::LSPV-4::SELECTOR-CAPACITY-BIT-LEDGER81